\documentclass{beamer}

\usetheme{Madrid}

% tiếng Việt cho pdflatex
\usepackage[utf8]{inputenc}
\usepackage[T5]{fontenc}
\usepackage[vietnamese]{babel}

\usepackage{amsmath}
\usepackage{tikz}

% màu pastel xanh
\definecolor{pastelblue}{RGB}{200,230,255}
\definecolor{softblue}{RGB}{150,200,235}
\definecolor{deepblue}{RGB}{60,120,170}

\setbeamercolor{structure}{fg=deepblue}
\setbeamercolor{frametitle}{bg=pastelblue, fg=black}

\title{Giới thiệu Lý thuyết Đồ thị}
\author{Nhóm 2}
\date{\today}

\begin{document}

%------------------------------------------------
\frame{\titlepage}

%------------------------------------------------
\begin{frame}{Mục lục}
\tableofcontents
\end{frame}

%------------------------------------------------
\section{Định nghĩa đồ thị}

\begin{frame}{Định nghĩa}

Đồ thị được ký hiệu:

\[
G = (V,E)
\]

Trong đó:

\begin{itemize}
\item $G$ : đồ thị (Graph)
\item $V$ : tập các đỉnh (Vertices)
\item $E$ : tập các cạnh (Edges)
\end{itemize}

\end{frame}

%------------------------------------------------
\begin{frame}{Ví dụ đồ thị}

\begin{center}

\begin{tikzpicture}

\node[circle,draw] (A) at (0,0) {A};
\node[circle,draw] (B) at (2,1) {B};
\node[circle,draw] (C) at (2,-1) {C};
\node[circle,draw] (D) at (4,0) {D};

\draw (A)--(B);
\draw (A)--(C);
\draw (B)--(D);
\draw (C)--(D);

\end{tikzpicture}

\end{center}

Ví dụ:

\[
V=\{A,B,C,D\}
\]

\[
E=\{AB,AC,BD,CD\}
\]

\end{frame}

%------------------------------------------------
\section{Đồ thị vô hướng}

\begin{frame}{Đồ thị vô hướng}

Trong đồ thị vô hướng:

\[
(u,v)=(v,u)
\]

Điều này nghĩa là cạnh không có hướng.

\end{frame}

%------------------------------------------------
\begin{frame}{Ví dụ đồ thị vô hướng}

\begin{center}

\begin{tikzpicture}

\node[circle,draw] (A) at (0,0) {A};
\node[circle,draw] (B) at (2,1) {B};
\node[circle,draw] (C) at (2,-1) {C};

\draw (A)--(B);
\draw (A)--(C);
\draw (B)--(C);

\end{tikzpicture}

\end{center}

Ví dụ:

\[
E=\{AB,AC,BC\}
\]

Ở đây:

\[
AB = BA
\]

\end{frame}

%------------------------------------------------
\section{Đồ thị có hướng}

\begin{frame}{Đồ thị có hướng}

Trong đồ thị có hướng:

\[
(u,v) \ne (v,u)
\]

Cạnh có hướng được gọi là \textbf{cung}.

\end{frame}

%------------------------------------------------
\begin{frame}{Ví dụ đồ thị có hướng}

\begin{center}

\begin{tikzpicture}[->]

\node[circle,draw] (A) at (0,0) {A};
\node[circle,draw] (B) at (2,1) {B};
\node[circle,draw] (C) at (2,-1) {C};

\draw (A)->(B);
\draw (B)->(C);
\draw (C)->(A);

\end{tikzpicture}

\end{center}

Ví dụ:

\[
E=\{(A,B),(B,C),(C,A)\}
\]

Ở đây:

\[
(A,B)\ne(B,A)
\]

\end{frame}

%------------------------------------------------
\section{Bậc của đỉnh}

\begin{frame}{Bậc của đỉnh}

Bậc của đỉnh là số cạnh nối với đỉnh đó.

Ký hiệu:

\[
deg(v)
\]

Trong đó:

\begin{itemize}
\item $deg$ : degree (bậc)
\item $v$ : đỉnh của đồ thị
\end{itemize}

\end{frame}

%------------------------------------------------
\begin{frame}{Ví dụ bậc của đỉnh}

\begin{center}

\begin{tikzpicture}

\node[circle,draw] (A) at (0,0) {A};
\node[circle,draw] (B) at (2,1) {B};
\node[circle,draw] (C) at (2,-1) {C};
\node[circle,draw] (D) at (4,0) {D};

\draw (A)--(B);
\draw (A)--(C);
\draw (A)--(D);
\draw (B)--(D);

\end{tikzpicture}

\end{center}

Bậc các đỉnh:

\[
deg(A)=3
\]

\[
deg(B)=2
\]

\[
deg(C)=1
\]

\[
deg(D)=2
\]

\end{frame}

%------------------------------------------------
\section{Định lý tổng bậc}

\begin{frame}{Định lý}

Trong đồ thị vô hướng:

\[
\sum deg(v)=2m
\]

Trong đó:

\begin{itemize}
\item $deg(v)$ : bậc của đỉnh $v$
\item $\sum$ : tổng tất cả các bậc
\item $m$ : số cạnh của đồ thị
\end{itemize}

\end{frame}

%------------------------------------------------
\begin{frame}{Ví dụ}

\begin{columns}

\column{0.45\textwidth}

\centering
\begin{tikzpicture}[scale=0.9]

\node[circle,draw] (A) at (0,0) {A};
\node[circle,draw] (B) at (2,1) {B};
\node[circle,draw] (C) at (2,-1) {C};

\draw (A)--(B);
\draw (A)--(C);
\draw (B)--(C);

\end{tikzpicture}

\column{0.55\textwidth}

Bậc các đỉnh:

\[
deg(A)=2,\quad deg(B)=2,\quad deg(C)=2
\]

Tổng bậc:

\[
2+2+2=6
\]

Số cạnh:

\[
m=3
\]

Ta có:

\[
2m=2\times3=6
\]

\end{columns}

\end{frame}
\end{document}
